\section{Cancellation conditions for an arbitrary number of stages}
\label{sec:general-order}

\subsection{A recursive representation of the ladder}

Consider a branch with $N\geq1$ stages, each containing a positive series resistance $r_j$ followed by a positive shunt capacitance $q_j$. The component normalization and frequency variable are those in Eq.~\eqref{eq:3}. The two-stage circuit in Figure~\ref{fig:circuit} is one instance of this class. With current directed from the network into the terminal load, the chain matrix of stage $j$ is
\begin{equation}\label{eq:general-matrix}
M_j(x)=\begin{bmatrix}1+r_jq_jx&r_j\\q_jx&1\end{bmatrix}.
\end{equation}
Let $(A_j,B_j)$ be the first row of $M_1\cdots M_j$. Appending the last stage gives
\begin{equation}\label{eq:general-forward}
B_j=r_jA_{j-1}+B_{j-1},\qquad
A_j=A_{j-1}+xq_jB_j,
\quad A_0=1,\ B_0=0.
\end{equation}
Induction in Eq.~\eqref{eq:general-forward} gives $\deg A_N=N$, $\deg B_N=N-1$ and $A_N(0)=1$. The leading coefficients are positive. Each stage matrix and the product have determinant one; the first-row polynomials have no nonconstant common factor.

The open-circuit transfer and source-shorted output impedance follow from the terminal equation. A common resistive load on each branch then gives the loaded transfer below, the same termination decomposition used in classical synthesis \citep{guillemin1949rcsynthesis}:
\begin{equation}\label{eq:general-transfer}
H_0=\frac1{A_N},\qquad
\frac{Z_0}{R_*}=\frac{B_N}{A_N},\qquad
H_g=\frac1{A_N+gB_N}.
\end{equation}
For $N=2$, $A_2=1+ax+bx^2$ and $B_2=c+dx$, recovering Eq.~\eqref{eq:5}. Retaining $B_N$ exposes the termination information absent from the open-circuit voltage transfer.

\subsection{Classical ladder extraction and coefficient coordinates}

With the source shorted, $B_N/A_N$ is the dimensionless driving-point impedance seen from the output. The driving-point impedance admits the following ladder continued fraction \citep{guillemin1949rcsynthesis,lee1978chain}
\[
z_j = \left[q_jx+\left(r_j+z_{j-1}\right)^{-1}\right]^{-1},
\qquad z_0=0,\quad z_N=B_N/A_N.
\]
The successive extraction below expresses continued-fraction ladder synthesis \citep{guillemin1949rcsynthesis,lee1978chain} in the $(A_j,B_j)$ coordinates, specialized to single-resistance/single-capacitance arms; recursive transfer calculations for general-arm ladders are classical as well \citep{fialkow1951rcladder}. We use the extraction to count the admissible component families under specified terminal constraints.

The last-stage values can be recovered from the leading coefficients. Writing $\operatorname{lc}(p)$ for the leading coefficient of polynomial $p$, define
\begin{equation}\label{eq:general-inverse}
\begin{aligned}
q_j&=\frac{\operatorname{lc}(A_j)}{\operatorname{lc}(B_j)},&
A_{j-1}&=A_j-xq_jB_j,\\
r_j&=\frac{\operatorname{lc}(B_j)}{\operatorname{lc}(A_{j-1})},&
B_{j-1}&=B_j-r_jA_{j-1}.
\end{aligned}
\end{equation}
For a positive ladder, these operations undo Eq.~\eqref{eq:general-forward} exactly. At each step, $\deg A_{j-1}=j-1$ and, when $j>1$, $\deg B_{j-1}=j-2$. The final pair is $(A_0,B_0)=(1,0)$. Each ratio is positive and uniquely specified.

Conversely, start with real polynomials satisfying $\deg A_N=N$, $\deg B_N=N-1$ and $A_N(0)=1$. If the reductions satisfy the stated highest-power constraints, every recovered component is positive, and the terminal pair is $(1,0)$, reversing the operations reconstructs the original polynomials from a positive ladder. The recovery conditions give a finite necessary and sufficient realization test for this topology. A failed recovery check rejects the candidate coefficient pair.

The forward map from the $2N$ positive components to the $2N$ coefficients of $A_N$ and $B_N$ is polynomial. The inverse map in Eq.~\eqref{eq:general-inverse} is rational with nonzero denominators at every positive realization. Strict positivity and the nonzero leading coefficients persist in a sufficiently small neighborhood. The maps are smooth inverses there, making the coefficient values valid local coordinates for the component family.

\subsection{The number of matching degrees of freedom}

Hold one reference branch fixed. Open-circuit matching fixes the $N$ nonconstant coefficients of $A_N$, leaving $N$ local degrees of freedom in the positive realization domain. Multiple positive realizations of one transfer function are classical; the present count concerns the component family that matches a fixed reference branch within this topology. Uniform impedance scaling changes $B_N$ by one scalar factor while preserving $A_N$. For $N\geq2$, that scaling follows one direction within a larger matching family.

At a specified $g>0$, raw matching, $H_{g,P}=H_{g,N}$, fixes all coefficients of $D_g=A_N+gB_N$. Termination previously entered realizability as a restriction on the attainable gain range \citep{fialkow1951rcladder,fialkow1952rcnetworks}; the count below instead concerns the family that matches a given reference branch. The constant coefficient fixes $B_N(0)$ through $D_g(0)=1+gB_N(0)$. The remaining $N$ coefficients fix the combinations of $A_N$ and $B_N$. The $N+1$ independent linear constraints leave $N-1$ local degrees of freedom, including zero for $N=1$.

Direct-current calibration instead compares $G_g=D_g(0)/D_g$. Fixing the normalized denominator imposes $N$ independent constraints and leaves $N$ local degrees of freedom. In coefficient coordinates the normalized coefficient of $x^N$ can be varied through $A_N$ independently of the lower coefficients. The constraint map therefore has rank $N$. The gain freedom used here has a classical counterpart in specified-gain synthesis, where the poles are fixed and the gain constant varies over a realizable interval \citep{fialkow1951rcladder,paige1960maximum,wang2017rcladder}. In the present matching problem that freedom is an ambiguity of the partner branch after calibration, with consequences for the uncalibrated differential output. For $N=1$, the positive family
\[
r>0,\qquad q=\frac{r+10}{11r}
\]
has the same normalized response at $g=0.1$ as $(r,q)=(1,1)$, although the raw direct-current gain varies with $r$.

For $N=2$, branch P = $(1,1,1,1)$ and the positive branch
\[
\left(2,2,\frac{217+\sqrt{6769}}{288},\frac{217-\sqrt{6769}}{480}\right)
\]
have identical normalized responses at $g=0.1$. The two raw denominators differ by the constant factor $7/6$, giving a raw CMRR of $20\log_{10}(13/2)=16.26$~dB at every frequency. Calibration therefore permits an additional matching direction that does not cancel the raw differential output.

\begin{table}[htbp]
\centering
\caption{Local matching dimensions under resistive terminations, for a fixed positive $N$-stage reference branch. Raw matching includes absolute gain; normalized matching compares $H_g/H_g(0)$.}
\label{tab:matching-conventions}
\begin{tabular}{lll}
\toprule
Termination information & Raw & Normalized \\
\midrule
Open circuit & $N$ & $N$ \\
One specified $g>0$ & $N-1$ & $N$ \\
Two distinct specified $g_1,g_2>0$ & $0$ & $0$ \\
\bottomrule
\end{tabular}
\end{table}

At two distinct positive conductances, put $U_i=1/H_{g_i}$. The two complete transfer functions determine
\begin{equation}\label{eq:general-two-load}
B_N=\frac{U_2-U_1}{g_2-g_1},\qquad
A_N=\frac{g_2U_1-g_1U_2}{g_2-g_1}.
\end{equation}
Two branches that agree at both loads consequently have identical $A_N$ and $B_N$. The stepwise uniqueness of the classical ladder extraction then requires identical corresponding components \citep{guillemin1949rcsynthesis}. The two-load reconstruction applies to every finite $N$ in the stated ladder class. It uses complete complex transfer functions, including absolute gain, at two loads.

Two normalized resistive-load responses also identify the components exactly. Write $D_i(g_k)=c_kD_{\rm ref}(g_k)$ for a candidate branch. The coefficient of $x^N$ is independent of $g_k$ and is nonzero; hence $c_1=c_2=c$. Subtracting the two load equations gives $B_i=cB_{\rm ref}$ and then $A_i=cA_{\rm ref}$. Since $A_i(0)=A_{\rm ref}(0)=1$, $c=1$. The normalized two-load proof uses exact complete responses. A stable noisy-data estimator additionally requires an uncertainty model.

\subsection{The load type also determines the matching family}

A known rational load can be written as a dimensionless admittance $g(x)=P(x)/Q(x)$, where $P$ and $Q$ are nonzero coprime real polynomials. The loaded response is defined provided $QA_N+PB_N$ is not identically zero. Equality of the two branch transfers gives
\begin{equation}\label{eq:rational-kernel}
H_g=\frac{Q}{QA_N+PB_N},\qquad
Q\Delta A+P\Delta B=0.
\end{equation}
Here $\Delta$ denotes branch P minus branch N. Coprimality implies $\Delta A=PS$ and $\Delta B=-QS$ for a polynomial $S$. The highest-power bounds and $\Delta A(0)=0$ restrict $S$ to
\begin{equation}\label{eq:rational-dimension}
\begin{aligned}
m&=\min\{N-\deg P,\ N-1-\deg Q\},\\
d_g&=\max\{0,m+1-\epsilon\},\qquad
\epsilon=\begin{cases}1,&P(0)\ne0,\\0,&P(0)=0.\end{cases}
\end{aligned}
\end{equation}
For $m<0$, the only allowable polynomial is $S=0$. Otherwise, $S$ has $m+1$ coefficients, and a nonzero $P(0)$ fixes the constant coefficient of $S$ to zero. Intersecting this linear family with the open positive realization domain preserves the local dimension $d_g$.

Table~\ref{tab:load-dimensions} applies Eq.~\eqref{eq:rational-dimension} to common loads. The open circuit, for which $P$ is identically zero, is listed separately. A pure capacitor changes the terminal capacitance and retains $N$ degrees of freedom, whereas a resistor also constrains the zero-frequency gain and reduces the count to $N-1$.
\begin{table}[htbp]
\centering
\caption{Local raw matching dimension for a fixed reference branch under one specified load, including absolute gain. All displayed load parameters are positive. Direct-current-normalized matching is treated separately in Table~\ref{tab:matching-conventions}.}
\label{tab:load-dimensions}
\begin{tabular}{lll}
\toprule
Load & Dimensionless admittance & Dimension \\
\midrule
Open circuit & $0$ & $N$ \\
Resistor & $g_0$ & $N-1$ \\
Capacitor & $\alpha x$ & $N$ \\
Inductor & $\beta/x$ & $\max(N-2,0)$ \\
Series RC & $\alpha x/(1+\gamma x)$ & $N-1$ \\
Parallel RC & $g_0+\alpha x$ & $N-1$ \\
\bottomrule
\end{tabular}
\end{table}

Two distinct known rational admittances also determine the component tuple under raw matching. For $g_k=P_k/Q_k$, multiply $Q_k\Delta A+P_k\Delta B=0$ by the other denominator and subtract. The result is $(P_1Q_2-P_2Q_1)\Delta B=0$. Nonidentical load functions mean $P_1Q_2-P_2Q_1\not\equiv0$, which gives $\Delta B=0$ and then $\Delta A=0$. Each known nonzero $Q_k$ permits cancellation as a rational function, including at isolated zeros by continuation. The loads in Table~\ref{tab:load-dimensions} are time-invariant rational admittances; a periodically switched input requires the separate sampling model introduced in Section~\ref{sec:dynamic}.

\subsection{Load variation and recovery conditioning}

For independent branches, let $D_i(g)=A_i+gB_i$. The exact common-mode-to-differential gain is
\begin{equation}\label{eq:general-leakage}
G_{DC}=-\frac{\Delta A+g\Delta B}{D_P(g)D_N(g)}.
\end{equation}
If the branches match at $g_0>0$, then $\Delta A=-g_0\Delta B$ and $\Delta B(0)=0$. Under a different common resistive load,
\begin{equation}\label{eq:general-load-shift}
G_{DC}(g)=-\frac{(g-g_0)\Delta B}{D_P(g)D_N(g)}.
\end{equation}
The zero-frequency leakage remains zero, while the remaining polynomial $\Delta B$ controls the frequency-dependent mismatch. The second-order load-shift formula is the case $\Delta B=x\Delta d$.

Exact identifiability and measurement conditioning require separate checks. Equation~\eqref{eq:general-two-load} divides by the load separation, and a small response perturbation gives $\delta(1/H)\simeq-\delta H/H^2$. Closely spaced loads and strongly attenuated responses amplify measurement errors. Experimental recovery therefore needs a stated frequency band and uncertainty model in addition to the algebraic uniqueness result.
For the dimensionless coefficient map $(r_1,\ldots,r_N,q_1,\ldots,q_N)\mapsto(A_{N,1},\ldots,A_{N,N},B_{N,0},\ldots,B_{N,N-1})$, the unscaled Jacobian at unit components has 2-norm condition numbers $15.6$, $1.08\times10^3$, $9.30\times10^4$, $8.67\times10^6$, $8.40\times10^8$ and $8.35\times10^{10}$ for $N=2,4,6,8,10,12$, respectively. The condition numbers concern the stated coefficient normalization and change under coordinate rescaling. They quantify local amplification in coefficient recovery; they are not a bound on the performance of a particular response-fitting estimator.
\FloatBarrier
